%==============================================================================
%  A parametric reformulation of the color structure tensor through virtual
%  illumination: equivalence and the limits of multiplicative extension to RGB-D
%
%  DRAFT SKELETON for submission to Pattern Recognition and Image Analysis (PRIA).
%
%  NOTE ON FORMAT: this file uses the standard `article` class so it compiles
%  anywhere with a basic TeX install (no journal class needed). Before
%  submission, port the body into the PRIA / Pleiades template
%  (pleiades.cls / spr-pria style). The theorems, equations, and tables below
%  are final and carry the actual measured numbers from the repository.
%
%  Placeholders to be written by the author are marked \todo{...} (red).
%  Draft prose already in place is plain text and may be edited freely.
%
%  Could not compile-test in the authoring environment (no TeX toolchain).
%  Build locally with:  latexmk -pdf article.tex   (or pdflatex x2).
%==============================================================================
\documentclass[11pt,a4paper]{article}

\usepackage[utf8]{inputenc}
\usepackage[T1]{fontenc}
\usepackage[margin=2.5cm]{geometry}
\usepackage{amsmath,amssymb,amsthm}
\usepackage{booktabs}
\usepackage{graphicx}
\usepackage{xcolor}
\usepackage[hidelinks]{hyperref}

% --- placeholder helpers -----------------------------------------------------
\newcommand{\todo}[1]{\textcolor{red}{[\textbf{TODO:} #1]}}
\newcommand{\note}[1]{\textcolor{blue}{\textit{#1}}}

% --- theorem environments ----------------------------------------------------
\theoremstyle{plain}
\newtheorem{theorem}{Theorem}
\newtheorem{proposition}{Proposition}
\newtheorem{corollary}{Corollary}
\theoremstyle{definition}
\newtheorem{definition}{Definition}

% --- notation shortcuts ------------------------------------------------------
\newcommand{\Lvec}{\mathbf{L}}
\newcommand{\grad}{\nabla}
\newcommand{\lmax}{\lambda_{\max}}
\newcommand{\lmin}{\lambda_{\min}}
\newcommand{\Mxy}{M_{XY}}
\newcommand{\transpose}{^{\mathsf T}}

\title{A Parametric Reformulation of the Color Structure Tensor
through Virtual Illumination: Equivalence and the Limits of
Multiplicative Extension to RGB-D}

% TODO: adapt author block to the PRIA template before submission.
% Per the author's preference, no e-mail / ORCID is embedded here.
\author{A.\,A.~Panchenko\\
  \small Ufa University of Science and Technology, Ufa, Russia
  % \and K.\,F.~Latypov  % TODO: confirm co-authorship / order
}
\date{Draft, \today}

\begin{document}
\maketitle

\begin{abstract}
\note{Draft abstract --- tighten before submission.}
Color edge detection that is sensitive to isoluminant transitions is
classically handled by the Di~Zenzo multichannel structure tensor.
We analyze a recently proposed ``virtual illumination'' edge operator,
which models the RGB channels as scalar surfaces lit from several virtual
directions, and show that it is \emph{not} an independent operator: with its
channel-asymmetric modulation disabled, its angularly-averaged response is a
strictly monotone function of the largest eigenvalue of the Di~Zenzo tensor,
so that threshold-binarized edge maps of the two methods coincide
(Spearman $\rho \approx 0.99$ on natural images, $1.00$ on synthetic edges).
We then characterize the single algebraically distinct ingredient of the
method --- a multiplicative height term --- and prove that it is structurally
unable to detect edges defined by the modulating channel alone, which closes
the RGB-D direction as a route to improving on the symmetric tensor. The
contribution is thus an explicit bridge between an intuitive physical model
and the established tensor formalism, a map of the admissible generalizations,
and an open, reproducible benchmark.
\todo{State the one-sentence headline result and quantitative highlights.}

\medskip
\noindent\textbf{Keywords:} color edge detection, structure tensor, Di~Zenzo
gradient, isoluminant edges, RGB-D, virtual illumination.
\end{abstract}

%==============================================================================
\section{Introduction}
\label{sec:intro}

\note{Draft --- expand.}
Edge detection on color images must, in principle, exploit chromatic
transitions that carry no luminance contrast (\emph{isoluminant} edges),
which grayscale operators such as Sobel, Prewitt, and Canny cannot see by
construction. A natural and intuitive way to reason about color gradients is
to picture each channel as a height surface and to ``illuminate'' it from
several virtual directions, accumulating the response. This paper studies such
a virtual-illumination operator and asks a precise question: \emph{is it a new
operator, or a re-derivation of an existing one?}

\todo{Motivate why the question matters: independent re-derivations are common;
establishing the exact relationship to prior art (Di~Zenzo) is the scientific
contribution, and it is non-obvious from the formula.}

\paragraph{Contributions.}
\begin{itemize}
  \item A parametric reformulation of the color structure tensor in which the
        lighting directions, their number, and a channel-asymmetric modulation
        appear as explicit, separately controllable degrees of freedom
        (Section~\ref{sec:decomp}).
  \item An equivalence theorem: with the modulation disabled, the operator
        reduces to the Di~Zenzo magnitude up to a strictly monotone map, hence
        to identical threshold-binarized output (Section~\ref{sec:equiv},
        Theorem~\ref{thm:equiv}).
  \item A structural-limit theorem for the multiplicative extension, and a
        resulting map of admissible RGB-D generalizations
        (Section~\ref{sec:rgbd}, Theorem~\ref{thm:mult}).
  \item An open, reproducible benchmark (synthetic, orientation, BSDS500,
        RGB-D) with all code released (Section~\ref{sec:exp}).
\end{itemize}

%==============================================================================
\section{Background: the Di~Zenzo structure tensor}
\label{sec:background}

\note{Draft --- expand with a short literature trace.}
For a multichannel image with channels $C_k$, Di~Zenzo~\cite{dizenzo1986}
defines the gradient through the structure tensor
\begin{equation}
\label{eq:dz-tensor}
M \;=\; \sum_{k} (\grad C_k)(\grad C_k)\transpose
\;=\;
\begin{pmatrix} E & F \\ F & G \end{pmatrix},
\quad
\begin{aligned}
E &= \textstyle\sum_k (C_{k,x})^2,\\
F &= \textstyle\sum_k C_{k,x} C_{k,y},\\
G &= \textstyle\sum_k (C_{k,y})^2,
\end{aligned}
\end{equation}
and takes the edge magnitude as $\sqrt{\lmax(M)}$, where $\lmax \ge \lmin \ge 0$
are the eigenvalues. The largest eigenvalue captures the maximal directional
variation across all channels jointly, so chromatic edges with no luminance
step are detected. \todo{Cite Cumani~\cite{cumani1991}, Sapiro~\cite{sapiro1996},
the color tensor of van de Weijer et al.~\cite{weijer2006}; note relation to
the grayscale structure tensor.}

%==============================================================================
\section{The virtual-illumination operator as a parametric decomposition}
\label{sec:decomp}

\note{Draft --- this section presents the method as the author's reformulation.}
Treat the RGB channels as three scalar surfaces. For an ordered channel triple
$(X,Y,Z)$ and a unit ``illumination'' direction $\Lvec=(d_x,d_y)\transpose$,
the operator forms the per-direction response
\begin{equation}
\label{eq:vl-response}
E_{\Lvec}
\;=\;
\sqrt{(\grad X \cdot \Lvec)^2 + (\grad Y \cdot \Lvec)^2}
\;\cdot\;
\underbrace{\Bigl(1 + \alpha\,\tfrac{C_Z - T}{255}\Bigr)}_{h(C_Z)},
\end{equation}
where $T$ is a global threshold (the $75$th percentile of the intensity
distribution) and $\alpha$ controls a channel-asymmetric \emph{height}
modulation. The response is averaged over a set of $N$ illumination directions
and, optionally, over the channel permutations, with a maximum fusion. The
free parameters --- the direction set $\{\Lvec_k\}$, their number $N$, the
modulation $h$ --- are exactly the degrees of freedom that, as we show next,
are folded into the single scalar $\lmax$ of the Di~Zenzo tensor.

\note{The key narrative for PRIA: this is a \emph{decomposition} that exposes
controllable knobs hidden inside the eigenvalue, not a new detector.}

%==============================================================================
\section{Equivalence to the Di~Zenzo tensor}
\label{sec:equiv}

The reduction rests on one identity. For any $u\in\mathbb{R}^2$ and unit
$\Lvec$,
\begin{equation}
\label{eq:identity}
(u\cdot\Lvec)^2 \;=\; \Lvec\transpose\,(u\,u\transpose)\,\Lvec .
\end{equation}
Summing \eqref{eq:identity} over the two gradient channels and setting
$\alpha=0$ (so $h\equiv1$),
\begin{equation}
\label{eq:quadform}
(\grad X\cdot\Lvec)^2 + (\grad Y\cdot\Lvec)^2
=
\Lvec\transpose
\underbrace{\bigl[(\grad X)(\grad X)\transpose + (\grad Y)(\grad Y)\transpose\bigr]}_{\Mxy}
\Lvec
= \Lvec\transpose \Mxy\, \Lvec ,
\end{equation}
where $\Mxy$ is precisely the two-channel Di~Zenzo tensor
\eqref{eq:dz-tensor}. Hence $E_{\Lvec}=\sqrt{\Lvec\transpose \Mxy\,\Lvec}$ is the
directional component of the structure tensor. In the eigenbasis of $\Mxy$,
with $\Lvec=(\cos\varphi,\sin\varphi)$,
\begin{equation}
\label{eq:radial}
E_\varphi \;=\; \sqrt{\lmax\cos^2\varphi + \lmin\sin^2\varphi}.
\end{equation}

\begin{theorem}[Equivalence]
\label{thm:equiv}
Let $\bar E = \frac{1}{2\pi}\int_0^{2\pi} E_\varphi\,d\varphi$ be the
angularly-averaged response \eqref{eq:radial}. Then
\begin{equation}
\label{eq:elliptic}
\bar E \;=\; \frac{2}{\pi}\,\sqrt{\lmax}\;
\mathrm{E}\!\left(\sqrt{1-\lmin/\lmax}\right),
\end{equation}
where $\mathrm{E}(\cdot)$ is the complete elliptic integral of the second
kind. For fixed anisotropy ratio $\lmin/\lmax$, $\bar E$ is a strictly
increasing function of $\sqrt{\lmax}$. Consequently any percentile threshold
selects the same pixels from $\bar E$ as from the Di~Zenzo magnitude
$\sqrt{\lmax}$; on rank-one edge structures ($\lmin=0$) the two binary edge
maps are identical.
\end{theorem}

\begin{proof}
\note{Draft --- formalize.}
Equation~\eqref{eq:elliptic} follows from \eqref{eq:radial} by the
substitution $k^2 = 1-\lmin/\lmax$ and the definition of $\mathrm{E}$.
Strict monotonicity in $\sqrt{\lmax}$ at fixed $k$ is immediate since
$\mathrm{E}(k)>0$. A strictly increasing transform preserves the order of
pixel values, hence preserves every percentile threshold set; for
$\lmin=0$, $\bar E=\frac{2}{\pi}\sqrt{\lmax}$ is exactly proportional to the
Di~Zenzo magnitude. \qed
\end{proof}

\begin{proposition}[Direction-set redundancy]
\label{prop:mode}
Because the response depends on $\Lvec$ only through the quadratic form
$\Lvec\transpose \Mxy\,\Lvec$, which is invariant under $\Lvec\mapsto-\Lvec$,
augmenting a direction set $\{\Lvec_1,\Lvec_2\}$ with its antipodes
$\{-\Lvec_1,-\Lvec_2\}$ leaves the averaged response unchanged. The
corresponding configurations (``mode~1'' and ``mode~2'') produce
pixel-identical maps.
\end{proposition}

\note{Empirically confirmed: \texttt{np.array\_equal} holds; see
Table~\ref{tab:corr} and the orientation study.}

%==============================================================================
\section{The multiplicative extension and its limit on RGB-D}
\label{sec:rgbd}

The single algebraically distinct ingredient of \eqref{eq:vl-response} is the
multiplicative height term $h(C_Z)$. With a fourth, semantically distinguished
channel $D$ (depth) placed in the role of $Z$, one might expect channel-
asymmetric modulation to help. We show the opposite, structurally.

\begin{theorem}[Structural limit of multiplicative modulation]
\label{thm:mult}
Let the full response be $E_{\Lvec}^{\mathrm{full}} = h(C_Z)\,
\sqrt{\Lvec\transpose \Mxy\,\Lvec}$ with $h$ bounded. If
$\grad X = \grad Y = 0$ at a pixel, then $E_{\Lvec}^{\mathrm{full}} = 0$ there
for every $\Lvec$, independently of $C_Z$ and $\grad C_Z$. Hence the operator
cannot respond to an edge defined solely by the $Z$ channel (e.g.\ a depth
step in a region of uniform color).
\end{theorem}

\begin{proof}
$\Mxy = 0$ when $\grad X=\grad Y=0$, so $\sqrt{\Lvec\transpose \Mxy\,\Lvec}=0$
and $h\cdot 0 = 0$ for any bounded $h$. \qed
\end{proof}

\begin{corollary}[Map of RGB-D generalizations]
\label{cor:map}
Any \emph{additive} incorporation of $D$, i.e.\ adding $(\grad D)(\grad D)
\transpose$ to the tensor, yields $\sqrt{\Lvec\transpose M_{RGBD}\,\Lvec}$ and
reduces by Theorem~\ref{thm:equiv} to multichannel Di~Zenzo. The only
algebraically distinct, \emph{multiplicative} extension cannot detect
$D$-only edges by Theorem~\ref{thm:mult}. Therefore, within this class, no
generalization both differs from Di~Zenzo and detects depth-defined edges;
doing so requires a mechanism nonlinear in $\grad D$.
\end{corollary}

%==============================================================================
\section{Experiments}
\label{sec:exp}

\note{All numbers below are the actual measured outputs of the released code
(\texttt{tests/}). Absolute timings are machine-dependent.}

\paragraph{Synthetic set.}
Table~\ref{tab:synth} reports F1 (2-pixel tolerance, $256\times256$).
Sobel/Prewitt/Di~Zenzo are binarized at their $95$th percentile.

\begin{table}[h]
\centering
\caption{Synthetic set, F1 at 2-pixel tolerance.}
\label{tab:synth}
\begin{tabular}{lcccc}
\toprule
Image & Sobel & Canny & DiZenzo & VL \\
\midrule
checkerboard                     & 0.000 & 0.945 & 0.739 & 0.933 \\
concentric\_circles              & 1.000 & 1.000 & 1.000 & 1.000 \\
color\_patches (equal bright.)   & 1.000 & 0.000 & 0.909 & 0.909 \\
color\_wheel                     & 0.510 & 0.471 & 0.340 & 0.887 \\
blurred\_disk                    & 1.000 & 1.000 & 1.000 & 1.000 \\
gray\_scale                      & 1.000 & 1.000 & 0.000 & 0.000 \\
\midrule
\textbf{Mean}                    & 0.752 & 0.736 & 0.665 & 0.788 \\
\bottomrule
\end{tabular}
\end{table}

\paragraph{Isoluminant orientation.}
On equal-brightness color stripes at $12$ orientations, both color methods
detect the edges (mean F1 $\approx 0.99$) while Canny scores $0.000$ at every
orientation. \todo{Insert the per-orientation curve / figure
(\texttt{orientation\_results.png}) or a compact summary table.}

\paragraph{Real images: BSDS500.}
Table~\ref{tab:bsds} reports the full $200$-image test split (union-of-
annotators ground truth, tolerance $0.0075\times\text{diagonal}$, fixed
thresholds; all detectors on the same images).

\begin{table}[h]
\centering
\caption{BSDS500 test split (200 images). Fixed thresholds (no ODS/OIS tuning).}
\label{tab:bsds}
\begin{tabular}{lccc}
\toprule
Method & F1 & Precision & Recall \\
\midrule
DiZenzo        & \textbf{0.575} & \textbf{0.654} & 0.531 \\
Canny          & 0.566 & 0.482 & \textbf{0.774} \\
VL (default)   & 0.552 & 0.630 & 0.509 \\
Prewitt        & 0.543 & 0.582 & 0.529 \\
Sobel          & 0.541 & 0.577 & 0.530 \\
\bottomrule
\end{tabular}
\end{table}

\note{All hand-crafted detectors sit far below modern learned methods
(BSDS ODS $\approx 0.83$--$0.84$); cite for context, do not compete on F1.}

\paragraph{Empirical equivalence.}
Table~\ref{tab:corr} reports the rank correlation between the continuous
response maps (Spearman $\rho$), confirming Theorem~\ref{thm:equiv}.

\begin{table}[h]
\centering
\caption{Map correlation (Spearman $\rho$). VL$(\alpha{=}0)$ vs.\ Di~Zenzo, and
the effect of the height term.}
\label{tab:corr}
\begin{tabular}{lcc}
\toprule
Image type & VL$(\alpha{=}0)\!\sim\!$DZ & VL$(\alpha{=}1)\!\sim\!$VL$(\alpha{=}0)$ \\
\midrule
Isoluminant patch        & 1.000 & 1.000 \\
Isoluminant stripes 30$^\circ$ & 0.996 & 0.996 \\
Real photos (BSDS, mean) & 0.99  & 0.97 \\
\bottomrule
\end{tabular}
\end{table}

\paragraph{Controlled RGB-D.}
Table~\ref{tab:rgbd} confirms Theorem~\ref{thm:mult}: VL-RGBD scores $0.000$ on
the depth-only edge, and $\alpha{=}1$ is indistinguishable from $\alpha{=}0$,
while a symmetric 4-channel Di~Zenzo handles every case.

\begin{table}[h]
\centering
\caption{Controlled synthetic RGB-D, F1 at 2-pixel tolerance.}
\label{tab:rgbd}
\begin{tabular}{lcccc}
\toprule
Case & DZ-RGB & DZ-RGBD & VL-RGBD ($\alpha{=}1$) & VL-RGBD ($\alpha{=}0$) \\
\midrule
color edge, flat depth      & 0.769 & 0.769 & 0.667 & 0.667 \\
\textbf{depth edge, flat color} & 0.000 & \textbf{0.769} & \textbf{0.000} & 0.000 \\
coincident color+depth      & 0.769 & 0.769 & 0.667 & 0.667 \\
orthogonal color/depth      & 0.722 & \textbf{0.910} & 0.636 & 0.636 \\
\midrule
\textbf{Mean}               & 0.565 & \textbf{0.790} & 0.490 & 0.490 \\
\bottomrule
\end{tabular}
\end{table}

%==============================================================================
\section{Related work}
\label{sec:related}
\todo{Position against: Di~Zenzo and multichannel gradients; Cumani's
multispectral edges; Sapiro's color/anisotropic work; the color structure
tensor (van de Weijer et al.); learned edge detection (HED and successors)
for context. Emphasize that the contribution is analytical (equivalence and
limits), not a new top-scoring detector.}

%==============================================================================
\section{Conclusion}
\label{sec:conclusion}
\note{Draft --- expand.}
We showed that an intuitive virtual-illumination color edge operator is, with
its channel-asymmetric modulation disabled, a re-derivation of the Di~Zenzo
structure tensor: the angularly-averaged response is a strictly monotone
function of $\sqrt{\lmax}$, giving identical threshold-binarized edges. The
only algebraically distinct ingredient, a multiplicative height term, is
structurally unable to detect edges defined by the modulating channel alone,
which closes the RGB-D direction within this class and leaves only nonlinear
mechanisms open. The practical by-product is an open, reproducible toolkit of
seven edge detectors and a benchmark.
\todo{One sentence on the broader takeaway: re-derivations are valuable when
their exact relationship to prior art is established and the boundary of the
new degrees of freedom is characterized.}

%==============================================================================
\section*{Reproducibility}
All code and data generators are released at
\url{https://github.com/Nervni-Sanya/Edge-detection-benchmark}.
The algebraic claims are mechanically reproduced by
\texttt{tests/correspondence\_analysis.py} (Theorem~\ref{thm:equiv}) and
\texttt{tests/rgbd\_synthetic\_test.py} (Theorem~\ref{thm:mult}); benchmark
tables by \texttt{tests/testing.py}. See \texttt{docs/equivalence.md} for the
full derivation.

%==============================================================================
\begin{thebibliography}{99}

\bibitem{dizenzo1986}
S.~Di~Zenzo,
``A note on the gradient of a multi-image,''
\emph{Computer Vision, Graphics, and Image Processing},
vol.~33, no.~1, pp.~116--125, 1986.

\bibitem{canny1986}
J.~Canny,
``A computational approach to edge detection,''
\emph{IEEE Trans. Pattern Anal. Mach. Intell.},
vol.~8, no.~6, pp.~679--698, 1986.

\bibitem{cumani1991}
A.~Cumani,
``Edge detection in multispectral images,''
\emph{CVGIP: Graphical Models and Image Processing},
vol.~53, no.~1, pp.~40--51, 1991.

\bibitem{sapiro1996}
G.~Sapiro and D.~L.~Ringach,
``Anisotropic diffusion of multivalued images with applications to color
filtering,''
\emph{IEEE Trans. Image Process.},
vol.~5, no.~11, pp.~1582--1586, 1996.

\bibitem{weijer2006}
J.~van de Weijer, T.~Gevers, and A.~W.~M.~Smeulders,
``Robust photometric invariant features from the color tensor,''
\emph{IEEE Trans. Image Process.},
vol.~15, no.~1, pp.~118--127, 2006.

\bibitem{arbelaez2011}
P.~Arbel\'aez, M.~Maire, C.~Fowlkes, and J.~Malik,
``Contour detection and hierarchical image segmentation,''
\emph{IEEE Trans. Pattern Anal. Mach. Intell.},
vol.~33, no.~5, pp.~898--916, 2011.

\bibitem{xie2015}
S.~Xie and Z.~Tu,
``Holistically-nested edge detection,''
in \emph{Proc. IEEE Int. Conf. Computer Vision (ICCV)}, 2015, pp.~1395--1403.

\bibitem{gonzalez2018}
R.~C.~Gonzalez and R.~E.~Woods,
\emph{Digital Image Processing}, 4th~ed.
Pearson, 2018.

% TODO: add a recent BSDS SOTA reference (e.g. EDTER / a 2024 baseline) for
% the learned-methods context, and any additional color-edge references.

\end{thebibliography}

\end{document}
